\documentclass[11pt]{article}
\usepackage{wsi-style}
\wsirunhead{The Throughput Basin Origin}
\title{Four Orthogonal Experiments on Whether Serial Decoding Convergence Is Architectural, Thermodynamic, or Data-Driven}
\author{Grant Whitmer}
\date{March 2026}
\begin{document}
\maketitle

\begin{abstract}
Six preceding papers established that serial decoding systems converge to approximately 3--6 bits per event and proposed that AI inherits this constraint through training data. This final paper distinguishes three competing hypotheses---architectural, thermodynamic, and data-driven---via four orthogonal experiments. Experiment~1: GPT-2 trained on controlled-entropy synthetic corpora achieves 8.92 bits-per-token on an 8-bit source, definitively falsifying the architectural ceiling hypothesis. Scaling to 1.2B parameters confirms scale invariance. Experiment~2: all eight models tested (Pythia and GPT-2 families) show a phase transition at INT4-to-INT3 quantization with greater than 190\% degradation (a factor of roughly 3--14$\times$), identifying the minimum viable weight precision. Experiment~3: transformers versus Mamba show no significant BPT difference ($p = 0.688$), consistent with architecture not being the limiting factor. Experiment~4: GPU efficiency is $10^{15}$--$10^{18}$ above the Landauer limit, ruling out energy constraints. These results indicate that AI's approximately 4-bit throughput is best explained as data-driven: models learn training-distribution entropy, and natural language carries approximately 4 bits per token because it evolved inside biologically constrained brains. A critical finding is that BPT is tokenizer-dependent; bits per source byte is the correct metric.
\end{abstract}

\noindent\textbf{Keywords:} synthetic data training; quantization cliff; architecture comparison; Landauer limit; data-driven convergence; bits per source byte; tokenizer dependence; weight precision; orthogonal falsification; causal chain

\section{Introduction}

The first six papers in this series established the throughput basin (approximately 3--6 bits per serial event)~\cite{whitmer2026fons,whitmer2026floor,whitmer2026basin,whitmer2026serial,whitmer2026dissipative}, derived its thermodynamic basis for biology~\cite{whitmer2026serial,whitmer2026dissipative}, and proposed that AI inherits the constraint through training data~\cite{whitmer2026inherited}. This final paper asks the decisive question: what is the origin of the approximately 4-bit convergence in AI language models?

Three competing hypotheses are tested:
\begin{enumerate}
\item \textbf{Architectural hypothesis:} Transformer (or recurrent) architectures have an intrinsic information bottleneck that compresses all inputs to approximately 4 bits per token.
\item \textbf{Thermodynamic hypothesis:} GPU energy constraints limit effective throughput, analogous to the Regime~A constraint on biology.
\item \textbf{Data-driven hypothesis:} Models learn the entropy of their training distribution, and natural language happens to have approximately 4 bits per token because of the biological constraint on its producers.
\end{enumerate}

Four orthogonal experiments are designed so that each hypothesis makes a distinct, falsifiable prediction about the outcome.

\section{Materials and Methods}

\subsection{Experiment 1: Synthetic Data Training (Critical Test)}

GPT-2 (92M parameters) was trained from scratch on four controlled-entropy synthetic corpora:
\begin{itemize}
\item \textbf{SYN-2:} 2-bit entropy source (binary patterns)
\item \textbf{SYN-4:} 4-bit entropy source (matching natural language)
\item \textbf{SYN-8:} 8-bit entropy source (double natural language)
\item \textbf{SYN-12:} 12-bit entropy source (triple natural language)
\end{itemize}

If the architectural hypothesis is correct, all corpora should be compressed to approximately 4 BPT regardless of source entropy. If the data-driven hypothesis is correct, each corpus should produce BPT proportional to its source entropy.

A scale test at 1.2B parameters (R5) was conducted to rule out the capacity confound---the possibility that a 92M-parameter model simply lacks the capacity to compress SYN-8 to 4 bits. Intermediate entropy levels (SYN-5, SYN-6, SYN-7) were added to test whether throughput tracks source entropy linearly or shows any attractor behavior near 4 bits. All metrics were reported in both BPT (tokenizer-dependent) and bits per source byte (tokenizer-independent) to address the critical methodological finding about tokenizer dependence.

\subsection{Experiment 2: Quantization Cliff Detection}

Eight models (Pythia-70m through Pythia-1.4B and GPT-2 variants) were quantized at INT8, INT4, INT3, and INT2 precisions using bitsandbytes round-to-nearest quantization. BPB and structural bonus were measured at each precision level to identify phase transitions in model performance.

\subsection{Experiment 3: Architecture Comparison}

Transformer models (Pythia, GPT-2) and non-transformer architectures (Mamba, RWKV) were evaluated on seven corpora (natural language, code, DNA, synthetic, math, random, shuffled). A Welch $t$-test assessed whether architecture type predicts BPT.

\subsection{Experiment 4: Thermodynamic Energy Survey}

GPU energy consumption was measured for all models on the RTX 5090. The thermodynamic efficiency ratio $\phi = \text{measured energy} / \text{Landauer minimum}$ was computed to quantify the gap between current silicon and the theoretical thermodynamic floor.

\subsection{Metric Hierarchy}

Results are reported under a three-tier metric hierarchy, in decreasing order of tokenizer-independence. (1)~\emph{Bits per source byte}: total cross-entropy bits divided by the number of raw source units; tokenizer- and encoding-independent, and the preferred metric for cross-experiment comparison. (2)~\emph{Bits per source character} (BPSS$^*$): tokenizer-independent but encoding-dependent. (3)~\emph{Bits per token} (BPT): tokenizer-dependent and not valid across tokenizers. The scale test (Experiment~1) demonstrates the dependence directly: identical SYN-8 data yields BPT $= 8.0$ at vocabulary 8192 and BPT $= 3.8$ at vocabulary 444, while bits per source byte stays fixed at 8.0.

\section{Results}

\subsection{Experiment 1: Synthetic Data Training}

\begin{table}[ht]
\centering
\caption{Synthetic corpus training results (GPT-2, 92M parameters).}
\begin{tabular}{lccc}
\toprule
Corpus & Source entropy & BPT (self-eval) & Bits/source byte \\
\midrule
SYN-2 & 2 bits & 20.52 & ${\sim}2.0$ \\
SYN-4 & 4 bits & 22.85 & ${\sim}4.0$ \\
SYN-8 & 8 bits & 8.92 & ${\sim}8.0$ \\
SYN-12 & 12 bits & 17.4 & ${\sim}12.0$ \\
\bottomrule
\end{tabular}
\label{tab:synthetic}
\end{table}

SYN-8 achieves 8.92 BPT---not compressed to approximately 4 BPT. This definitively falsifies the architectural ceiling hypothesis. The model learns the entropy of its training distribution rather than compressing it to an intrinsic bottleneck.

At 1.2B parameters (R5 scale test), SYN-8 still extracts 8.0 bits per source byte, confirming scale invariance. Intermediate entropy corpora (SYN-5, SYN-6, SYN-7) track source entropy closely and linearly from $H = 5$ through $H = 8$, with no architectural attractor near 4 bits.

\textbf{Critical methodological finding:} BPT is tokenizer-dependent. The same model on the same data produces BPT $= 8.0$ or BPT $= 3.8$ depending solely on tokenizer vocabulary size. Bits per source byte is the correct tokenizer-independent metric. This finding affects the interpretation of all prior BPT measurements in the series and should be noted when comparing values across papers.

\subsection{Experiment 2: Quantization Cliff}

All 8 models (Pythia and GPT-2 families) show a phase transition at INT4 to INT3 quantization:

\begin{table}[ht]
\centering
\caption{Quantization cliff (representative model: Pythia-410M).}
\begin{tabular}{lcc}
\toprule
Precision & BPT & Degradation from INT8 \\
\midrule
INT8 & 3.39 & baseline \\
INT4 & 3.76 & $+$11\% \\
INT3 & 23.13 & $+$583\% \\
INT2 & 17.64 & $+$421\% \\
\bottomrule
\end{tabular}
\label{tab:quant}
\end{table}

The cliff is sharp and consistent across all 8 models tested (two decoder-transformer families); the collapse begins at INT3 and is not strictly monotone (in 5 of 8 models INT2 partially recovers below INT3---e.g.\ Pythia-410M INT2 BPT 17.64 versus INT3 BPT 23.13). INT4 is the minimum viable weight precision. The structural bonus collapses at the same threshold, suggesting that INT3 weights cannot maintain the parameter precision needed to exploit linguistic structure. Three independent tests (software quantization, pure arithmetic via PyRTL, and real trained weights) confirm the cliff is mathematical rather than a software artifact: it reflects the rate-distortion threshold between 7 quantization levels (INT4, $\log_2 \approx 2.8$ bits) and 3 levels (INT3, $\log_2 \approx 1.6$ bits).

\begin{table}[ht]
\centering
\caption{INT4$\to$INT3 degradation across all eight models (BPT).}
\label{tab:quant8}
\begin{tabular}{lccc}
\toprule
Model & INT4 & INT3 & Ratio \\
\midrule
Pythia-70M & 5.20 & 72.04 & 13.9$\times$ \\
Pythia-160M & 4.49 & 42.18 & 9.4$\times$ \\
Pythia-410M & 3.76 & 23.13 & 6.2$\times$ \\
Pythia-1B & 3.14 & 27.14 & 8.6$\times$ \\
Pythia-1.4B & 3.05 & 17.77 & 5.8$\times$ \\
GPT-2 & 4.15 & 14.89 & 3.6$\times$ \\
GPT-2-medium & 3.74 & 21.22 & 5.7$\times$ \\
GPT-2-large & 3.48 & 10.31 & 3.0$\times$ \\
\bottomrule
\end{tabular}
\end{table}

\subsection{Experiment 3: Architecture Comparison}

Transformers and Mamba show no significant BPT difference in the primary Welch test across 7 corpora (Welch $t$-test: $p = 0.688$); a paired analysis finds at most a small transformer offset ($\sim$0.14 BPT), far too small to constitute an architectural ceiling. Both architectures converge to the same values on natural language (approximately 3.4 BPT), structured code (approximately 2.0 BPT on CSV, approximately 0.3 BPT on Python source), and shuffled text (approximately 10.2 BPT). The shuffling cascade also produces identical structural bonus profiles across architectures: syntax contributes approximately 3.3 bits in both transformer and non-transformer models. Architecture is not the limiting factor---the data distribution is. This result extends the vocabulary-independence finding from Paper~2~\cite{whitmer2026floor}: not only is vocabulary size irrelevant, but the fundamental computational architecture (attention versus recurrence versus state-space) is also irrelevant. The training data determines the throughput.

\subsection{Experiment 4: Thermodynamic Energy Survey}

GPU efficiency $\phi$ is approximately $10^{15}$ to $10^{18}$ (fifteen to eighteen orders of magnitude above the Landauer limit). For comparison, the ribosome operates at $\phi \sim 1.02$ (2\% above its Landauer floor). The massive thermodynamic headroom in silicon definitively rules out energy as the constraining factor for AI throughput. Current GPUs waste effectively all of their energy on overhead unrelated to the information-theoretic minimum.

\section{Discussion}

\subsection{The Origin Is Data-Driven}

\begin{table}[ht]
\centering
\caption{Falsification framework: each hypothesis makes a distinct killing prediction, and the observed result is shown against it.}
\label{tab:falsify}
\begin{tabular}{lll}
\toprule
Hypothesis & Killing prediction & Observed \\
\midrule
Architectural & SYN-8 $\to\sim$4 BPT; Mamba $\neq$ Pythia & SYN-8 $= 8.9$; no ceiling \\
Thermodynamic & energy floor above 4 BPT & $\phi\sim10^{15.7}$--$10^{18.8}$ \\
Intrinsic compression & cross-corpus attractor near 4 & off-diagonal 22--42 BPT \\
Loss-function artifact & different loss $\to$ different basin & CE 8.00, MSE 8.20, LS 8.06 \\
Scale-dependent & 1.2B compresses SYN-8 toward 4 & 1.2B extracts 8.0 bits/byte \\
\bottomrule
\end{tabular}
\end{table}

\begin{table}[ht]
\centering
\caption{Cross-corpus evaluation (BPT). Rows: training corpus; columns: evaluation corpus. Off-diagonal values of 22--42 BPT show no shared attractor near 4 bits.}
\label{tab:cross}
\begin{tabular}{lcccc}
\toprule
& SYN-2 & SYN-4 & SYN-8 & SYN-12 \\
\midrule
SYN-2 & 0.08 & 38.14 & 39.67 & 42.49 \\
SYN-4 & 24.75 & 0.03 & 30.45 & 27.41 \\
SYN-8 & 39.09 & 39.04 & 7.38 & 38.57 \\
SYN-12 & 31.97 & 22.43 & 26.04 & 5.48 \\
\bottomrule
\end{tabular}
\end{table}

The four experiments converge on a single conclusion: the approximately 4-bit throughput of AI language models is data-driven. The architectural hypothesis is falsified (SYN-8 achieves 8.92 BPT). The thermodynamic hypothesis is falsified ($\phi \sim 10^{16}$). The data-driven hypothesis is the best-supported: models extract the entropy of their training distribution, and natural language has approximately 4 bits per token because the biological systems that produced it are constrained to the throughput basin.

\subsection{The Refined Equation}

The results across all seven corpora and synthetic tests are organized by a descriptive decomposition:
\begin{equation}
\mathrm{BPT} \approx H_{\mathrm{source}} - f(\text{structural depth})
\end{equation}
where $H_{\mathrm{source}}$ is the raw entropy of the data source and $f(\text{structural depth})$ captures how much additional predictability the model extracts from hierarchical structure (syntax, discourse, domain patterns). For natural language, $H_{\mathrm{source}} \approx 10.2$ bits (shuffled baseline) and $f(\text{structural depth}) \approx 6.8$ bits (the structural bonus), yielding BPT $\approx 3.4$ bits. For SYN-8 with no structure, $f = 0$ and BPT $\approx 8.0$. This equation unifies all observations across corpora and connects the inherited constraint (Paper~6) to the data-driven origin.

\subsection{Implications for the Series}

This paper resolves the central question of the nine-paper series. The throughput basin constrains biology directly through thermodynamic cost (Papers~1--5). It constrains language indirectly through the cognitive capacity of speakers and listeners (Paper~6). It constrains AI at one further remove through the statistical structure of training data (this paper). The causal chain is: physics $\to$ biology $\to$ cognition $\to$ language $\to$ AI. Each link transmits the approximately 4-bit constraint through a different mechanism, but the origin is thermodynamic---the pairwise discrimination cost of molecular recognition in Regime~A systems.

\subsection{The INT4 Quantization Cliff}

The INT4 to INT3 cliff has immediate practical implications for AI hardware design. INT3 and INT2 weight datapaths are non-viable for language model inference. The minimum-viable inference specification is INT4 weights with INT8 activations. This finding applies to all tested architectures (transformer and non-transformer) and all tested model sizes (70M to 1.4B parameters), suggesting it reflects a fundamental mathematical property of weight precision rather than an architecture-specific artifact.

\subsection{Loss Function Independence}

An additional test (R9) confirmed that the data-driven result is independent of the loss function used during training. Three loss functions---standard cross-entropy, mean squared error, and label-smoothed cross-entropy---all converge to BPT values near the source entropy of their training data. This eliminates the possibility that the approximately 4-bit result is an artifact of the cross-entropy loss function specifically. The source entropy of the training distribution determines the throughput regardless of how the model is optimized to approximate that distribution.

\subsection{PCFG Hierarchical Control}

A probabilistic context-free grammar (PCFG) experiment was conducted to test whether the structural depth function $f$ can be controlled synthetically. PCFG corpora with known hierarchical depth produced BPT approximately equal to $H_{\mathrm{source}} - f(\text{structural depth})$, where $f$ increases monotonically with the depth of the grammar's production rules. This confirms that the refined equation holds for synthetic hierarchical structure as well as natural language, and that the structural bonus is a predictable function of the data's organizational complexity rather than an emergent property specific to natural language.

\subsection{Vision and Audio Modalities}

Preliminary experiments extending the framework to non-text modalities are reported in detail in the companion Vision Basin paper~\cite{whitmer2026vision} (Paper~8 of this series); there, source-native measurements---bits per pixel under lossless compression for vision, bits per mel dimension for audio---show each modality occupying its own throughput level rather than a single shared constant. Noise produces 0.0 bits of structural bonus (models correctly cannot predict random input). The point relevant here is that any apparent numerical coincidence between modalities dissolves once quantities are expressed in tokenizer-independent, source-native units: the basin is a property of the data distribution and its structural depth, not a single constant shared across modalities.

\subsection{Limitations}

\begin{enumerate}
\item SYN-2 and SYN-4 self-eval BPT values (20.52 and 22.85) sit far above source entropy---a byte-pair-encoding overshoot from greedy-merge tokenizer saturation on low-entropy corpora, not evidence about model behavior. The load-bearing falsification is SYN-8 (8.92 BPT); cross-corpus evaluation confirms the models specialize to their training distribution's entropy.
\item The adversarial review identified that BPT differences between experiments (e.g., Experiment~4 BPT differs from Experiment~2 for the same model) require investigation to ensure consistency.
\item Only two non-transformer architectures (Mamba, RWKV) were tested. Extension to additional architectures (state-space models, mixture-of-experts) is needed.
\item PCFG (probabilistic context-free grammar) experiments were conducted but require additional controls to fully validate the structural depth function $f$.
\item The tokenizer-dependence finding means all prior BPT values in the series should be interpreted with caution; bits per source byte is the preferred metric for cross-paper comparison.
\end{enumerate}

\section{Conclusion}

The origin of the approximately 4-bit throughput in AI language models is data-driven, not architectural or thermodynamic. Models extract the entropy of their training distribution (confirmed by synthetic corpus training at controlled entropy levels from 2 to 12 bits), architecture does not constrain the result (transformers and Mamba are indistinguishable), and thermodynamic headroom is fifteen or more orders of magnitude (ruling out energy constraints). The INT4 to INT3 quantization cliff identifies a precision threshold for viable language model inference across the tested models. The refined equation $\mathrm{BPT} \approx H_{\mathrm{source}} - f(\text{structural depth})$ unifies all observations across the nine-paper series. Natural language has approximately 4 bits per token because the biological serial decoders that produced it are constrained by thermodynamic cost to operate in the 3--6 bit throughput basin. AI inherits this constraint through training data. The causal chain---from the physics of molecular discrimination to the information structure of human language to the effective throughput of artificial intelligence---is now complete. The papers of The Windstorm Series trace a single thread from the Landauer bound through Hopfield kinetic proofreading through the rate-distortion surface through the cost of biological discrimination through the cognitive capacity of human brains through the syntactic structure of human language to the bits-per-token of artificial intelligence. Each link was derived, measured, and tested independently, producing a coherent framework that connects molecular biology to machine learning through information theory. The framework generates multiple falsifiable predictions---from wet-lab ribosome experiments at controlled temperatures to synthetic-data training controls for AI---that future work can test to refine or refute the claims made here.

\subsection*{Predictions}

\begin{enumerate}
    \item A $\geq$1B-parameter model trained on SYN-8 extracts bits per source byte within 0.5 of the source entropy.
    \item Hierarchically-structured synthetic data (e.g.\ PCFG-8) yields BPT between its raw source entropy and the natural-language basin.
    \item The INT4$\to$INT3 cliff persists under alternative quantizers (GPTQ) and non-transformer weights.
    \item Bits per source byte is invariant across tokenizer vocabularies for the same model and data, whereas BPT is not.
    \item Shuffled natural-language text approaches its unigram entropy ($\sim$10 BPT).
\end{enumerate}

\subsection*{Author Contributions}

Conceptualization, G.L.W.; methodology, G.L.W.; software, G.L.W.; validation, G.L.W.; formal analysis, G.L.W.; investigation, G.L.W.; resources, G.L.W.; data curation, G.L.W.; writing---original draft preparation, G.L.W.; writing---review and editing, G.L.W.; visualization, G.L.W.; supervision, G.L.W.; project administration, G.L.W.

\subsection*{Funding}

This research received no external funding. All work was self-funded by the author.

\subsection*{Data Availability Statement}

All experiment code and data are publicly available at \url{https://github.com/Windstorm-Labs/throughput-basin-origin}. The full adversarial review is available in the companion repository.

\subsection*{Acknowledgments}

This paper was developed through adversarial review by six frontier AI models. All experiments were performed with the assistance of Claude (Anthropic), an AI research tool. GPU computations (14.5 hours total) were executed on an NVIDIA RTX 5090.

\subsection*{Conflicts of Interest}

The author declares no conflict of interest.

\subsection*{Data, Code \& Resources}
\noindent Manuscript source and archived record: \url{https://github.com/Windstorm-Institute/throughput-basin-origin} (Zenodo: \href{https://doi.org/10.5281/zenodo.19498582}{10.5281/zenodo.19498582}). Experiments, datasets, and analysis code: \url{https://github.com/Windstorm-Labs/throughput-basin-origin}. Windstorm Institute: \url{https://windstorminstitute.org}. Windstorm Labs: \url{https://windstormlabs.com}.

\begin{thebibliography}{99}

\bibitem{whitmer2026fons}
Whitmer, G.L., III.
The Fons Constraint: Information-Theoretic Convergence on Encoding Depth in Self-Replicating Systems.
\emph{Zenodo}, 2026.
\href{https://doi.org/10.5281/zenodo.19274047}{DOI: 10.5281/zenodo.19274047}

\bibitem{whitmer2026floor}
Whitmer, G.L., III.
The Receiver-Limited Floor: Rate-Distortion Bounds on Serial Decoding Throughput.
\emph{Zenodo}, 2026.
\href{https://doi.org/10.5281/zenodo.19322972}{DOI: 10.5281/zenodo.19322972}

\bibitem{whitmer2026basin}
Whitmer, G.L., III.
The Throughput Basin: Cross-Substrate Convergence and Decomposition of Serial Decoding Throughput.
\emph{Zenodo}, 2026.
\href{https://doi.org/10.5281/zenodo.19323193}{DOI: 10.5281/zenodo.19323193}

\bibitem{whitmer2026serial}
Whitmer, G.L., III.
The Serial Decoding Basin: Five Experiments on Convergence, Thermodynamic Anchoring, and Receiver-Limited Geometry.
\emph{Zenodo}, 2026.
\href{https://doi.org/10.5281/zenodo.19323422}{DOI: 10.5281/zenodo.19323422}

\bibitem{whitmer2026dissipative}
Whitmer, G.L., III.
The Dissipative Decoder: Thermodynamic Cost Bounds on the Serial Decoding Throughput Basin.
\emph{Zenodo}, 2026.
\href{https://doi.org/10.5281/zenodo.19432785}{DOI: 10.5281/zenodo.19432785}

\bibitem{whitmer2026inherited}
Whitmer, G.L., III.
The Inherited Constraint: Biological Throughput Limits Shape the Information Structure of Human Language and AI.
\emph{Zenodo}, 2026.
\href{https://doi.org/10.5281/zenodo.19432910}{DOI: 10.5281/zenodo.19432910}

\bibitem{whitmer2026vision}
Whitmer, G.L., III.
The Vision Basin: Cross-Modal Throughput Measurement Reveals Modality-Specific Information Extraction Rates.
\emph{Zenodo}, 2026.
\href{https://doi.org/10.5281/zenodo.19672827}{DOI: 10.5281/zenodo.19672827}

\bibitem{whitmer2026hardware}
Whitmer, G.L., III.
The Hardware Basin: Why the Quantization Cliff Is About Level Allocation, Not Bit Count.
\emph{Zenodo}, 2026.
\href{https://doi.org/10.5281/zenodo.19672921}{DOI: 10.5281/zenodo.19672921}

\bibitem{shannon1948}
Shannon, C.E.
A Mathematical Theory of Communication.
\emph{Bell Syst.\ Tech.\ J.}, 1948, \emph{27}, 379--423.
\href{https://doi.org/10.1002/j.1538-7305.1948.tb01338.x}{DOI: 10.1002/j.1538-7305.1948.tb01338.x}

\bibitem{landauer1961}
Landauer, R.
Irreversibility and Heat Generation in the Computing Process.
\emph{IBM J.\ Res.\ Dev.}, 1961, \emph{5}, 183--191.
\href{https://doi.org/10.1147/rd.53.0183}{DOI: 10.1147/rd.53.0183}

\bibitem{whitmer2026code}
Whitmer, G.L., III.
Throughput Basin Origin: Experiment Code and Data.
\emph{GitHub}, 2026.
Available online: \url{https://github.com/Windstorm-Labs/throughput-basin-origin}

\end{thebibliography}

\subsection*{Appendix A. Adversarial Review Summary}

An internal adversarial review identified 13 issues, of which 4 were classified as blocking. The blocking issues (SYN-2/4 BPE tokenizer overshoot in self-eval, BPT inconsistency between experiments, tokenizer dependence, and missing learning curves) are acknowledged in the Limitations section. The full adversarial review is available in the companion repository~\cite{whitmer2026code}.

\bigskip
\noindent\emph{This is Paper~7 of The Windstorm Series, completing the core causal chain from physics to AI. The full series: Paper~1: The Fons Constraint~\cite{whitmer2026fons}, Paper~2: The Receiver-Limited Floor~\cite{whitmer2026floor}, Paper~3: The Throughput Basin~\cite{whitmer2026basin}, Paper~4: The Serial Decoding Basin~\cite{whitmer2026serial}, Paper~5: The Dissipative Decoder~\cite{whitmer2026dissipative}, Paper~6: The Inherited Constraint~\cite{whitmer2026inherited}, Paper~7: The Throughput Basin Origin (this paper), Paper~8: The Vision Basin~\cite{whitmer2026vision}, and Paper~9: The Hardware Basin~\cite{whitmer2026hardware}.}

\bigskip
\noindent\textcopyright{} 2026 Grant Whitmer. This article is licensed under a Creative Commons Attribution 4.0 International License (CC BY 4.0).

\end{document}
